% Section IV: Method -- the network, its loss, training and regularization,
% the hybrid inputs, tuning and the final-network rule, the ablation study and
% the metrics. Settings: src/indoorloc/config.py (DL_FIXED, DL_GRID, DL_EPOCHS,
% DL_VARIANTS); chosen setting: results/phase2/tuning/best_params.json.
% Weight counts: notebook 04, section 6 (indoorloc.dl.MultiTaskMLP).
\section{Method}\label{sec:method}

\begin{figure}[t]
\centering
\begin{tikzpicture}[
  font=\footnotesize, >=latex, node distance=3mm and 4mm,
  box/.style={draw=black!60, rounded corners=2pt, align=center, inner sep=3pt,
              minimum height=7mm, fill=white},
  input/.style={box, minimum width=17mm},
  hybrid/.style={box, minimum width=17mm, fill=blue!8, draw=blue!60!black},
  layer/.style={box, minimum width=9mm, minimum height=22mm, fill=black!5},
  head/.style={box, minimum width=15mm}]
\node[input] (rssi) {Scaled RSSI\\(464)};
\node[input, below=of rssi] (heard) {Heard\\indicators (464)};
\node[hybrid, below=of heard] (phase1) {Phase~1 features\\(3 + 8)};
\node[layer, right=6mm of heard] (h1) {512};
\node[layer, right=of h1] (h2) {256};
\node[layer, right=of h2] (h3) {128};
\node[head, right=6mm of h3] (floor) {Floor\\(5 classes)};
\node[head, above=of floor] (building) {Building\\(3 classes)};
\node[head, below=of floor] (position) {Position\\$(x, y)$};
\foreach \n in {rssi, heard, phase1} \draw[->] (\n.east) -- (h1.west);
\draw[->] (h1) -- (h2);
\draw[->] (h2) -- (h3);
\foreach \n in {building, floor, position} \draw[->] (h3.east) -- (\n.west);
\node[above=1mm of h2, align=center] {Shared trunk\\(ReLU, dropout)};
\node[above=1mm of building, align=center] {Heads};
\end{tikzpicture}
\caption{Our hybrid multi-task network. The shared trunk turns a scan's input
columns into 128 numbers that all three heads read. The blue block is the
hybrid part: the Phase~1 features (3 signal statistics and 8 AP-affinity
columns).}
\label{fig:network}
\end{figure}

\subsection{The Multi-Task Network}\label{sec:network}
Our network is a multilayer perceptron (MLP) with hard parameter
sharing~\cite{caruana1997multitask,ruder2017overview}: three fully connected
hidden layers (512, 256 and 128 units) are shared by all tasks, and three small
output layers, the \emph{heads}, read the last one (Fig.~\ref{fig:network}).
The widths are not tuned, because every setting costs a five-fold CV of 200
epochs; we use a funnel that halves the width at each layer, of the size used by
other networks on this dataset (shared layers of 512--512
in~\cite{laska2022multitask}, 1024--512--256 in~\cite{li2026hybrid}). A hidden
layer computes
\begin{equation}
  \mathbf{h}_{l}=\operatorname{ReLU}(W_{l}\,\mathbf{h}_{l-1}+\mathbf{b}_{l}),
  \qquad \operatorname{ReLU}(a)=\max(0,a),
  \label{eq:layer}
\end{equation}
with weights $W_l$ and biases $\mathbf{b}_l$; $\mathbf{h}_0$ is the input
vector. The building and floor heads give one score $z_c$ per class, and the
softmax $p_c=e^{z_c}/\sum_k e^{z_k}$ turns the scores into probabilities; the
prediction is the class with the highest one. The position head gives the two
standardized coordinates.

\subsection{Loss}\label{sec:loss}
For a mini-batch of $N$ scans, the building head uses the cross-entropy
$L_B=-\frac1N\sum_i\log p_{i,y_i}$, the mean of $-\log$ of the probability
given to the true class, the floor head the same loss $L_F$ for floors, and the
position head the mean squared error of the
standardized coordinates, $L_P=\frac1{2N}\sum_i\lVert\hat{\mathbf{u}}_i-
\mathbf{u}_i\rVert^2$. Because $x$ and $y$ are standardized with their own
training standard deviations (124.5 and 67.4~m), $L_P$ counts a 1~m error in $y$
about 3.4 times as much as a 1~m error in $x$, while the 2D error counts both
equally. The network minimizes one weighted sum,
\begin{equation}
  L = L_B + L_F + \lambda\,L_P .
  \label{eq:loss}
\end{equation}
Because the trunk is shared, its gradient is the sum of the three tasks'
gradients, $\nabla_{\theta}L=\nabla_{\theta}L_B+\nabla_{\theta}L_F+
\lambda\nabla_{\theta}L_P$: every task shapes the shared representation,
which is how multi-task learning shares information between
tasks~\cite{caruana1997multitask}. The weight $\lambda$ balances the tasks; a
position error of 10~m is only about 0.1 in standardized units, so with
$\lambda=1$ the classification losses would dominate. Results are known to be
sensitive to these weights~\cite{kendall2018multi}, so $\lambda$ is tuned.

\subsection{Training and Regularization}\label{sec:training}
We implement the networks in PyTorch~\cite{paszke2019pytorch} and train them
with AdamW~\cite{kingma2015adam,loshchilov2019decoupled}, mini-batches
of 256 scans and a learning rate of $10^{-3}$ that falls to 0 along a cosine
curve over 200 epochs (passes over the training scans), so that the last
steps are small. Three things stop the network from memorizing the
training scans: \emph{dropout}~\cite{srivastava2014dropout} sets each hidden
unit to 0 with probability $p=0.2$ during training (only then),
\emph{weight decay} (0.01, the default of PyTorch's AdamW) shrinks every
weight a little at each step, separately from the Adam step as AdamW
proposes~\cite{loshchilov2019decoupled}: PyTorch computes
$\theta\leftarrow\theta-\eta\,(0.01\,\theta+\ldots)$ with the current learning
rate $\eta$; and \emph{hidden WAPs}: in every batch, each heard WAP of a scan is
hidden with probability 0.1 (its scaled RSSI and indicator are set to 0; the 11
Phase~1 columns are left as they are). A new phone, or the same phone on another
day, misses some WAPs, so the network should not depend on any single one. This
is a form of input dropout~\cite{srivastava2014dropout}, but without dropout's
rescaling: a hidden WAP looks exactly like a WAP that is not heard.

\subsection{Two Hybrids of Phase~1 and Phase~2}\label{sec:hybrid}
The raw RSSI columns say how strongly each WAP is heard, but not where the WAP
is; a network must learn that from the training scans by itself. Our hybrid
network also receives the Phase~1 features (Section~\ref{sec:inputs}): the
AP-affinity columns are a location hint (likely building, floor and position)
computed from where each WAP is heard in the training scans. This is
classical, hand-made feature engineering feeding a deep network. Wide-and-deep
models~\cite{cheng2016wide} also join hand-made features with a deep network,
but through a separate linear part added at the output; ours are plain extra
inputs of the shared trunk. The AP-affinity profiles are
learned again in every fold, so the hint never uses held-out labels.

The second hybrid combines the models at the output. The Phase~1 Random
Forest and the hybrid MLP each give class probabilities for building and
floor and a position; the \emph{blend} averages them,
\begin{equation}
  \hat{\mathbf{q}}=w\,\mathbf{q}_{\mathrm{RF}}+(1-w)\,\mathbf{q}_{\mathrm{NN}},
  \label{eq:blend}
\end{equation}
with one forest weight $w\in\{0,0.1,\dots,1\}$ per task ($w=0$: the network
alone, $w=1$: the forest alone) and the most probable class for building and
floor. The weights are chosen by user-group CV: in every fold both models are
trained on the other folds and predict the held-out scans, and each task keeps
the weight with the best CV score, as each task kept its own setting in
Phase~1 (on a tie, the smaller weight). The network's out-of-fold outputs use
the ablation seed (Section~\ref{sec:tuning}), not the runs that won the tuning,
so the network alone ($w=0$) is the ablation's full network. Trees and networks
are biased towards different kinds of functions~\cite{grinsztajn2022tree}, in
Phase~1 no model was best at every task, and on tabular data an ensemble of
networks and XGBoost beat XGBoost alone~\cite{shwartzziv2022tabular}, so an
average can beat both models.

\subsection{Tuning, Final Network and Ablation}\label{sec:tuning}
We tune the one setting that only a multi-task network has, the position
weight $\lambda\in\{30,100,300\}$, because multi-task results depend on the task
weights~\cite{kendall2018multi}; the other settings are fixed in advance
(Section~\ref{sec:training}). Every value is scored on the same five user-group
folds as Phase~1, and the one with the lowest CV mean 2D error is kept. Before using the test set we fixed the final network: the \emph{hybrid
MLP} of Fig.~\ref{fig:network}, the shared design of the brief. Two other
networks are tested next to it with the same settings: the \emph{MLP}
(multi-task, Wi-Fi inputs only) and \emph{hybrid single-task MLPs} (three
separate networks, one per task, with the hybrid inputs). The ablation study removes one part of the
hybrid MLP at a time (multi-task sharing, Phase~1 features, heard indicators,
hidden WAPs, dropout and weight decay) and scores each variant with the same
folds and settings. All ablation networks use a new random seed, so that the
full network is not the very run that won the tuning (which would favor it).
Each fold has a different difficulty, so we also count in how many of the five
folds a variant is worse. To see how much a CV score depends on chance, we
cross-validate the full network with two more seeds (a seed sets the starting
weights, the dropout masks, the order of the batches and the hidden WAPs). Finally, we train the final network on all
training scans with five seeds and test each, to see the spread of the test
score.

\subsection{Metrics}\label{sec:metrics}
We use exactly the Phase~1 metrics and code: building and floor accuracy, the
mean absolute error (MAE) of $x$ and $y$, the mean and median 2D error
(straight-line distance between predicted and true position), the error
breakdown into three cases (building wrong; building right but floor wrong;
both right), floor confusion matrices, error maps (in notebook~05) and the
split into test scans from phones seen in training and from new phones.
